\documentclass{article}

% ================= Packages =================
\usepackage{xcolor}
\usepackage{tikz}
\usetikzlibrary{calc,positioning}
\usepackage{setspace}
\usepackage{geometry}
\usepackage{graphicx}
\usepackage{listings}
\usepackage{xepersian}

% ================= Page Layout =================
\linespread{1.3}
\geometry{margin=2.5cm}
\setlength{\parindent}{0.7cm}
\setlength{\hfuzz}{9pt}

% ================= Colors =================
\definecolor{main_blue}{RGB}{0,70,140}
\definecolor{light_blue}{RGB}{0,120,200}
\definecolor{light_gray}{RGB}{112,146,190}
\definecolor{code_bg}{RGB}{245,245,250}
\definecolor{code_kw}{RGB}{0,70,140}
\definecolor{code_str}{RGB}{180,60,0}
\definecolor{code_cm}{RGB}{100,130,100}

% ================= Fonts =================
\settextfont[
Scale=1.3,
Path=../../report/template/fonts/,
Extension=.TTF,
BoldFont=B-NAZANIN-Bold
]{B-NAZANIN}

\setlatintextfont[
Scale=1.1,
Path=../../report/template/fonts/,
Extension=.TTF,
BoldFont=timesbd,
ItalicFont=timesi
]{Times}

% ================= Code Listing =================
\lstset{
  language=Python,
  backgroundcolor=\color{code_bg},
  basicstyle=\ttfamily\small,
  keywordstyle=\color{code_kw}\bfseries,
  stringstyle=\color{code_str},
  commentstyle=\color{code_cm}\itshape,
  numbers=left,
  numberstyle=\tiny\color{light_gray},
  frame=single,
  framerule=0.5pt,
  rulecolor=\color{light_gray},
  tabsize=4,
  breaklines=true,
  showstringspaces=false,
  xleftmargin=0.5cm,
}

% ================= Helpers =================
\renewcommand{\labelitemi}{\textcolor{main_blue}{\large$\bullet$}}
\newcommand{\pellipsis}{\ensuremath{\ldots}}
\newcommand{\pmdash}{\lr{---}}

% ================= Custom Section =================
\newcommand{\mysection}[2]{
\par\vspace{1cm}
\noindent
\begingroup
\def\mysectionNumW{1.2cm}
\def\mysectionH{1.2cm}
\def\mysectionR{3pt}
\def\mysectionGap{2pt}
\begin{tikzpicture}
    \coordinate (rightedge) at (\linewidth,0);
    \node[rectangle, fill=light_blue, text=white, anchor=east,
    minimum width=\mysectionNumW, minimum height=\mysectionH,
    rounded corners=\mysectionR, align=center, inner sep=2pt, outer sep=0pt,
    font=\Large\bfseries] (num) at (rightedge) {#1};
    \node[rectangle, fill=main_blue, text=white, anchor=east,
    minimum height=\mysectionH, rounded corners=\mysectionR,
    align=right, inner sep=6pt, outer sep=0pt,
    font=\Large\bfseries] (title) at ([xshift=-\mysectionGap]num.west) {#2};
    \draw[light_gray, line width=3pt, line cap=round]
    ($(0,0 |- num.north)+(0.5\pgflinewidth,0)$) -- ($(rightedge |- num.north)+(-0.5\pgflinewidth,0)$);
\end{tikzpicture}
\endgroup
\par\vspace{0.5cm}
}

\newcommand{\mysubsection}[1]{
\par\vspace{0.5cm}
\noindent\textcolor{main_blue}{\large\textbf{#1}}
\par\vspace{0.2cm}
}

% ================= Info Box =================
\newcommand{\infobox}[1]{
\vspace{0.3cm}
\noindent
\begin{tikzpicture}
  \node[rectangle, fill=code_bg, draw=light_gray, rounded corners=3pt,
        inner sep=8pt, text width=\dimexpr\textwidth-1.2cm\relax] {#1};
\end{tikzpicture}
\vspace{0.3cm}
}

% ================= Document =================
\begin{document}

\begin{figure}[h]
\centering
\includegraphics[width=4cm]{../../report/template/AUT.png}
\end{figure}

\vspace{0.3cm}
\begin{center}
{\LARGE\textbf{راهنمای کامل کد پروژه}}\\[0.4cm]
{\large تشخیص خرابی ماشین‌آلات با \lr{Transformer} روی \lr{FPGA}}\\[0.3cm]
{\large مهرداد شیخ عباسی}\\[0.2cm]
{\normalsize دانشگاه امیرکبیر \pmdash{} دانشکده مهندسی کامپیوتر}
\end{center}

\vspace{0.3cm}
\noindent{\color{light_gray}\rule{\linewidth}{0.5mm}}
\vspace{0.3cm}

\noindent
این سند راهنمای خواندنِ کامل کدِ پروژه است \pmdash{} از بارگذاریِ داده تا خروجیِ \lr{FPGA}.
هر بخش یک فایل از پایپلاین را توضیح می‌دهد: هدفش چیست، کدِ کلیدی‌اش چه می‌کند،
و چطور به بقیه وصل است.

% ==============================================================
\mysection{۱}{\lr{src/config.py} \pmdash{} تنظیمات مرکزی}
% ==============================================================

\mysubsection{هدف}

\lr{config.py} تنها جایی است که اعدادِ مهم پروژه نگه داشته می‌شوند.
هیچ اسکریپتی نباید عددِ \lr{hard-code} داشته باشد \pmdash{} همه از اینجا \lr{import} می‌کنند.
این یعنی برای تغییرِ معماری یا تنظیمات فقط یک فایل را باید عوض کرد.

\mysubsection{دسته‌بندیِ ثابت‌ها}

\begin{itemize}
\item \textbf{مسیرها:} \lr{PROJECT\_ROOT}، \lr{DATA\_RAW\_DIR}، \lr{RUNS\_DIR} \pmdash{} پروژه را جابه‌جا کنی مسیرها خودکار درست می‌شوند.

\item \textbf{داده:} \lr{WINDOW\_SIZE = 1024}، \lr{WINDOW\_STRIDE = 1024}،
\lr{INPUT\_CHANNELS = 8}، \lr{NORMALIZATION\_SCOPE = "per\_window"}.

\item \textbf{کلاس‌ها:} \lr{NUM\_CLASSES = 4}، \lr{LABEL\_MAP} (نگاشتِ نام به عدد):
\end{itemize}

\begin{latin}
\begin{lstlisting}
LABEL_MAP = {
    "normal":       0,
    "imbalance":    1,
    "misalignment": 2,
    "bearing":      3,
}
\end{lstlisting}
\end{latin}

\begin{itemize}
\item \textbf{معماری \lr{Transformer}:}
\lr{TRANSFORMER\_D\_MODEL = 32}،
\lr{TRANSFORMER\_NHEAD = 4}،
\lr{TRANSFORMER\_NUM\_LAYERS = 1}،
\lr{TRANSFORMER\_DIM\_FEEDFORWARD = 64}.
این چهار عدد مدلِ \lr{F\_1layer} را کاملاً تعریف می‌کنند (\lr{9,492} پارامتر).

\item \textbf{آموزش:} \lr{NUM\_EPOCHS}، \lr{BATCH\_SIZE}، \lr{LEARNING\_RATE}،
\lr{RANDOM\_SEED = 42}،\\
\lr{EARLY\_STOPPING\_PATIENCE}.
\end{itemize}

\mysubsection{نکته کلیدی: چرا \lr{d\_model = 32}؟}

\lr{d\_model} تعدادِ ویژگیِ هر توکن در طولِ مدل است. در جست‌وجوی معماری،
۱۱ پیکربندی از \lr{d=8} تا \lr{d=64} تست شد. \lr{F\_1layer} با \lr{d=32} روی زانوی منحنیِ
دقت\pmdash{}اندازه قرار دارد: کوچک‌ترین مدلی که تقریباً ۹۸ درصد دقت می‌دهد.

\infobox{
\textbf{قاعده:} \lr{d\_model \% nhead == 0} باید برقرار باشد.
برای ما: $32 \div 4 = 8$ (\lr{head\_dim = 8}، کمینه قابل‌قبول).
}

% ==============================================================
\mysection{۲}{\lr{src/data/load\_data.py} \pmdash{} کشفِ فایل‌ها و برچسب‌گذاری}
% ==============================================================

\mysubsection{هدف}

این ماژول پوشه \lr{MAFAULDA} را اسکن می‌کند، همه فایل‌های \lr{CSV} را پیدا می‌کند،
و از روی نامِ مسیر به هر فایل یک برچسبِ کلاس می‌دهد.

\mysubsection{تابعِ اصلی: \lr{build\_file\_index}}

\begin{latin}
\begin{lstlisting}
def build_file_index(data_dir) -> pd.DataFrame:
    # returns: columns = [file_path, label, split_hint]
\end{lstlisting}
\end{latin}

خروجی یک \lr{DataFrame} است که هر ردیفش یک فایل \lr{CSV} از \lr{MAFAULDA} است.

\mysubsection{برچسب‌گذاری با \lr{infer\_label\_from\_path}}

\begin{latin}
\begin{lstlisting}
def infer_label_from_path(path: Path) -> str | None:
    rel = path.relative_to(DATA_RAW_DIR)  # only path under data/raw
    rel_str = rel.as_posix().lower()
    if "normal"       in rel_str: return "normal"
    if "imbalance"    in rel_str: return "imbalance"
    if "misalignment" in rel_str: return "misalignment"
    if "bearing"      in rel_str: return "bearing"
    return None
\end{lstlisting}
\end{latin}

\infobox{
\textbf{نکته مهم:} تابع فقط روی بخشِ زیرِ \lr{DATA\_RAW\_DIR} جست‌وجو می‌کند،
نه کلِ مسیرِ مطلق. اگر اسمِ پوشه پروژه حاوی کلیدواژه‌هایی مثل \lr{"normal"} باشد،
برچسب‌گذاری اشتباه نمی‌شود. (باگی که پیدا و اصلاح شد.)
}

\mysubsection{جداسازیِ فایل‌ها: \lr{split\_file\_index}}

\begin{latin}
\begin{lstlisting}
train_df, val_df, test_df = split_file_index(index_df)
# default ratio: 70% / 10% / 20%
\end{lstlisting}
\end{latin}

جداسازی در سطحِ \textbf{فایل} انجام می‌شود، نه پنجره. این از نشتِ داده (\lr{data leakage}) جلوگیری می‌کند:
هیچ پنجره‌ای از یک فایل نمی‌تواند هم در \lr{train} و هم در \lr{test} باشد.

% ==============================================================
\mysection{۳}{\lr{make\_dataloaders.py} \pmdash{} پنجره‌سازی و \lr{DataLoader}}
% ==============================================================

\mysubsection{هدف}

از فایل‌های \lr{CSV} به تنسورهای \lr{PyTorch} می‌رسیم.
مسیر: بارگذاریِ \lr{CSV} $\leftarrow$ نرمال‌سازی $\leftarrow$ پنجره‌بندی $\leftarrow$ \lr{DataLoader}.

\mysubsection{پنجره‌بندی}

هر فایل \lr{CSV} یک سیگنالِ زمانی است. آن را به پنجره‌های ثابت‌اندازه تقسیم می‌کنیم:

\begin{latin}
\begin{lstlisting}
WINDOW_SIZE   = 1024  # samples per window
WINDOW_STRIDE = 1024  # step size (no overlap)
# one window shape: (1024, 8)
\end{lstlisting}
\end{latin}

\mysubsection{نرمال‌سازی: \lr{per-window z-score}}

\begin{latin}
\begin{lstlisting}
# per-window independently:
mu  = window.mean()
std = window.std()
normalized = (window - mu) / (std + eps)
\end{lstlisting}
\end{latin}

\begin{itemize}
\item هر پنجره مستقل نرمال می‌شود: میانگین صفر، واریانس یک.
\item برای \lr{FPGA} ایده‌آل است: پنجره اطلاعاتِ کافی دارد و نیاز به آمارِ سراسری نیست.
\item دامنه مطلق (شدتِ خرابی) از دست می‌رود، ولی برای تشخیصِ \textbf{نوعِ} خرابی کافی است.
\end{itemize}

\mysubsection{تابعِ اصلی: \lr{build\_dataloaders}}

\begin{latin}
\begin{lstlisting}
train_loader, val_loader, test_loader, stats = build_dataloaders(
    batch_size=64,
    window_size=1024,
    stride=1024,
    normalization_scope="per_window",
)
# each batch: X shape (64, 1024, 8), y shape (64,)
\end{lstlisting}
\end{latin}

% ==============================================================
\mysection{۴}{\lr{src/models/transformer\_baseline.py} \pmdash{} معماریِ مدل}
% ==============================================================

\mysubsection{مسیرِ داده در مدل}

\begin{latin}
\begin{lstlisting}
(batch, 1024, 8)
    -> Linear(8->32)         # input_projection
    -> PositionalEncoding    # + positional info
    -> EncoderBlock          # self-attention + FFN  (1 block)
    -> mean(dim=1)           # (batch, 32)
    -> Linear(32->16)->GELU->Linear(16->4)   # classifier
(batch, 4)                   # logit for 4 classes
\end{lstlisting}
\end{latin}

\mysubsection{کلاسِ اول: \lr{PositionalEncoding}}

مشکل: \lr{attention} ترتیبِ توکن‌ها را نمی‌داند.
راه‌حل: به هر جایگاه یک بردارِ ثابتِ یکتا اضافه می‌کنیم.

\begin{latin}
\begin{lstlisting}
pe = zeros(1024, 32)
pe[:, 0::2] = sin(position * div_term)  # even columns
pe[:, 1::2] = cos(position * div_term)  # odd columns

# in forward:
x = x + pe[:, :seq_len, :]
\end{lstlisting}
\end{latin}

\begin{itemize}
\item \lr{div\_term}: اعدادی از ۱.۰ تا ۰.۰۰۰۲ که هر بُعد را با آهنگِ متفاوت نوسان می‌دهد.
\item هر جایگاه یک ترکیبِ \lr{(sin, cos)} یکتا دارد \pmdash{} مثلِ شمارشِ دودویی با موجِ نرم.
\item \lr{register\_buffer}: جدول \lr{train} نمی‌شود، ثابت است.
\end{itemize}

\mysubsection{کلاسِ دوم: \lr{TransformerEncoderBlock}}

\textbf{۱) \lr{Self-Attention}:} هر توکن به همه توکن‌ها نگاه می‌کند.

\begin{latin}
\begin{lstlisting}
attn_out, _ = self.self_attention(x, x, x)  # Q=K=V=x
x = LayerNorm(x + Dropout(attn_out))        # Add & Norm
\end{lstlisting}
\end{latin}

سه \lr{x} به این معنی نیست که مقادیر یکسانند \pmdash{} داخلِ \lr{MultiheadAttention}
با وزن‌های جداگانه $W_Q, W_K, W_V$ ضرب می‌شوند.
\lr{nhead=4}: توجه ۴ بار موازی با دیدگاه‌های مختلف (\lr{head\_dim = 32/4 = 8}).

\textbf{۲) \lr{Feedforward}:} روی هر توکن جداگانه:

\begin{latin}
\begin{lstlisting}
ff_out = Linear(32->64) -> GELU -> Linear(64->32)
x = LayerNorm(x + Dropout(ff_out))   # Add & Norm
\end{lstlisting}
\end{latin}

\mysubsection{کلاسِ سوم: \lr{TransformerBaseline}}

\begin{latin}
\begin{lstlisting}
def forward(self, x):
    x = self.input_projection(x)      # (B,1024,8) -> (B,1024,32)
    x = self.positional_encoding(x)   # (B,1024,32)
    for layer in self.encoder_layers:
        x = layer(x)                  # (B,1024,32)
    x = x.mean(dim=1)                 # (B,32)  Global Avg Pool
    return self.classifier(x)         # (B,4)
\end{lstlisting}
\end{latin}

\infobox{
\lr{mean(dim=1)}: میانگین روی ۱۰۲۴ توکن \pmdash{} از \lr{(batch, 1024, 32)} به \lr{(batch, 32)}.
این \lr{Global Average Pooling} نام دارد: ساده‌ترین روش برای خلاصه کلِ پنجره.
}

% ==============================================================
\mysection{۵}{\lr{src/train\_transformer\_baseline.py} \pmdash{} حلقه آموزش}
% ==============================================================

\mysubsection{حلقه اصلی}

\begin{latin}
\begin{lstlisting}
for epoch in range(num_epochs):
    model.train()
    for X_batch, y_batch in train_loader:
        optimizer.zero_grad()          # clear previous gradients
        logits = model(X_batch)        # forward pass
        loss = criterion(logits, y_batch)
        loss.backward()                # compute gradients
        optimizer.step()               # update weights

    val_loss, val_acc = evaluate(model, val_loader, ...)

    if val_loss < best_val_loss:
        best_model_state = deepcopy(model.state_dict())
        torch.save({...}, checkpoint_path)  # save
    else:
        epochs_without_improvement += 1
        if epochs_without_improvement >= patience:
            break   # early stopping
\end{lstlisting}
\end{latin}

\begin{itemize}
\item \textbf{\lr{optimizer}:} \lr{Adam} با \lr{lr=1e-3}
\item \textbf{\lr{criterion}:} \lr{CrossEntropyLoss} (اگر \lr{imbalance > 1.5} با وزنِ کلاس)
\item \textbf{بهترین مدل:} بر اساسِ \lr{val\_loss} انتخاب می‌شود
\item \textbf{\lr{checkpoint}:} شاملِ وزن‌ها \texttt{+} معماری \texttt{+} تنظیمات
\end{itemize}

\mysubsection{توابعِ مشترک (\lr{src/utils/training\_utils.py})}

\begin{latin}
\begin{lstlisting}
set_seed(42)               # fix all random sources
evaluate(...)              # loss and accuracy on DataLoader
evaluate_with_details(...) # + confusion matrix
build_loss_function(...)   # CrossEntropy with/without class weights
\end{lstlisting}
\end{latin}

% ==============================================================
\mysection{۶}{\lr{src/evaluate/evaluate\_model.py} \pmdash{} ارزیابی}
% ==============================================================

\mysubsection{لودِ مدل از چک‌پوینت}

\begin{latin}
\begin{lstlisting}
checkpoint = torch.load(path, map_location=device)
model_config = checkpoint.get("model_config")
model = TransformerBaseline(
    d_model=model_config.get("d_model", TRANSFORMER_D_MODEL),
    # missing value falls back to config.py default
)
model.load_state_dict(checkpoint["model_state_dict"])
\end{lstlisting}
\end{latin}

معماری از \textbf{داخلِ چک‌پوینت} خوانده می‌شود نه \lr{config.py} فعلی.
\lr{map\_location=device}: اگر مدل روی \lr{GPU} ذخیره شده ولی الان \lr{CPU} داری، درست کار می‌کند.

% ==============================================================
\mysection{۷}{\lr{src/quantize/} \pmdash{} کوانتیزاسیونِ \lr{int10}}
% ==============================================================

\mysubsection{هدف}

مدلِ \lr{float32} برای \lr{FPGA} مناسب نیست. راه‌حل: همه اعداد را به \lr{int10} تبدیل کنیم.

\begin{center}
\begin{tabular}{|c|c|c|}
\hline
\textbf{تعدادِ بیت} & \textbf{دقت} & \textbf{افتِ دقت} \\
\hline
\lr{float32} & \lr{97.55\%} & \pmdash{} \\
\lr{int16}   & \lr{97.55\%} & \lr{0.00\%} \\
\lr{int10}   & \lr{97.11\%} & \lr{0.44\%} \\
\lr{int8}    & کمتر & بیشتر \\
\hline
\end{tabular}
\end{center}

\lr{int10} نقطه بهینه است: افتِ دقت ناچیز، منابعِ \lr{FPGA} بسیار کمتر.

\mysubsection{فرآیند}

\begin{latin}
\begin{lstlisting}
# step 1: calibrate on train subset
scales = calibrate(model, X_train)

# step 2: integer inference
# float_input -> /scale_in -> int
# MatMul with int weights -> *scale_out -> next layer input
\end{lstlisting}
\end{latin}

\mysubsection{فایل‌های فعال}

\begin{itemize}
\item \lr{integer\_ops.py}: عملیاتِ پایه (\lr{LayerNorm} صحیح، \lr{GELU} با \lr{LUT}، \lr{Softmax} صحیح)
\item \lr{integer\_model.py}: کلاسِ \lr{IntegerTransformer} با \lr{numpy}
\item \lr{evaluate\_integer.py}: کالیبراسیون + مقایسه \lr{float} در برابر \lr{integer}
\item \lr{evaluate\_integer\_bits.py}: جست‌وجوی بهترین تعدادِ بیت (یافت: \lr{int10})
\item \lr{evaluate\_int10\_metrics.py}: متریکِ نهاییِ \lr{int10}
\item \lr{export\_hls\_f\_int10.py}: \lr{export} وزن‌ها و \lr{scale}ها برای \lr{HLS C++}
\end{itemize}

% ==============================================================
\mysection{۸}{\lr{hls\_f\_int10/} \pmdash{} هسته \lr{C++} برای \lr{FPGA}}
% ==============================================================

\mysubsection{هدف}

پیاده‌سازیِ مدلِ \lr{int10} در \lr{C++} برای سنتزِ \lr{HLS} (\lr{Vitis HLS}).

\begin{latin}
\begin{lstlisting}
hls_f_int10/
  core.cpp / core.hpp    # inference logic
  data/                  # weights and scales (from export_hls_f_int10.py)
  golden/                # reference input/output for testbench
  onchip_normalization_reference.py  # NumPy reference for on-chip z-score
\end{lstlisting}
\end{latin}

تأییدِ درستی قبل از سنتز: \lr{python -m hls\_f\_int10.verify\_oracle}
نتیجه: \lr{PASS 16/16}.

% ==============================================================
\mysection{۹}{\lr{pynq/, host/} \pmdash{} درایور و رابطِ کاربری}
% ==============================================================

\begin{itemize}
\item \lr{pynq/pynq\_driver.py}: اجرای مدل روی برد \lr{EBAZ4205 (Zynq-7010)}
\item \lr{pynq/numpy\_backend.py}: شبیه‌سازی بدونِ برد (برای تست)
\item \lr{host/uart\_monitor.py}: دریافتِ پیامِ \lr{serial} از \lr{FPGA}
\item \lr{host/display\_ui.py}: نمایشِ نتیجه روی صفحه کامپیوتر
\item پروتکل: \lr{\texttt{FRAME=<id>,CLASS=<n>,LABEL=<name>,CONF=<c>}}
\end{itemize}

% ==============================================================
\mysection{۱۰}{\rl{تصویرِ کلی از ورودی تا خروجی}}
% ==============================================================

\begin{latin}
\begin{lstlisting}
MAFAULDA CSV files (8 channels, 50 kHz)
    |
    v  load_data.py       : file discovery + labeling
    v  make_dataloaders.py: windowing (1024) + per-window z-score
    v  transformer_baseline.py: 8->32->attention->32->4 classes
    v  (train, best checkpoint saved)
    |
    v  IntegerTransformer : calibrate scales + int10 inference
    v  export_hls_f_int10 : weights + scales -> data/ files
    v  HLS C++ core       : synthesized for EBAZ4205
    v  PYNQ driver        : feed window -> get class
    v  UART -> host PC    : display_ui.py -> show result
\end{lstlisting}
\end{latin}

\end{document}
